% O5. Existence at the prior 1/2: part (d) of the full statement of Proposition 3.
% Placement: Appendix A.4, after the proof of parts (a) to (c) of the full statement of Proposition 3, as the proof
% of part (d). The named block repeats part (d) for reading; in the appendix the statement appears once, inside the
% full statement (statements/B.tex), and the proof follows the proofs of parts (a) to (c).
% Source: scan/reduction_A3_A5.md section 3(d), read again here as a full proof; PLAN.md Part 5, O5.
% Cites: Lemma A.2 (lem:prices), Lemma A.3 (lem:trading) and Lemmas A.7 to A.9 (lem:a4-high, lem:a4-low,
% lem:a4-product) in the form of statements/B.tex, where Lemma A.7 writes (Q) as 0<s_+<lambda with mu_+(s_+)<tau,
% so that it covers tau>=1.

\begin{named}{Proposition~\ref{prop:coexistence}(d), existence at the prior \(1/2\)}
Let the classes satisfy the class ranking, let \(F\) be continuous with \(\DT>0\) and \(g_H>g_L\), let the prior be \(1/2\) and
\(\tau>1/2\), including \(\tau\ge1\), and let an indifferent investor mix over its best orders. Under either challenger
rule an equilibrium exists against \(F\). If \(\rho\DT\le2k\), the quiet market is one. If \(\rho\DT>2k\), let
\(t=\min\{\lambda,\,b\,(\rho\DT/k-2)\}\). If \(\mu_+(t)<\tau\), both classes trade at the rate \(t\) and entry is
\(\rho\) after every order. If \(\mu_+(t)\ge\tau\), the low class sells at the rate \(t\), class \(H\) buys at the rate
\(\min\{\lambda,\,b\,(\DT/k-2)\}\ge t\), a purchase recruits a high-cost challenger, and entry after a sale and after
the zero order is \(\rho\).
\end{named}

\begin{proof}[Proof of Proposition~\ref{prop:coexistence}(d)]
Fix \(F\) and the slot, and let the prior be \(1/2\). Since \(\pi<\tau\), Lemmas~\ref{lem:a4-high} to~\ref{lem:a4-product}
apply against \(F\). The proof has five steps.

\emph{Step 1: the frame.} At the prior \(1/2\),
\[
1-\mu_+(s)=\frac{b}{2b+s}=\mu_-(s)\qquad\text{for every }s\in[0,\lambda],
\]
and in particular \(1-M=m=b/(2b+\lambda)\). Each tuple below has no wrong-signed order, rates with \(s_+\ge s_-\),
entry \(\rho\) after a sale and after the zero order, a stated entry \(e(1)\) after a purchase, the prices
\(P(x)=t_0+e(x)\,w(\mu(x))\), and the challenger's belief \(\mu(x)\) at the price \(P(x)\). By
Lemma~\ref{lem:a4-product}, such a tuple is an equilibrium under a pair of tie rules when \((s_+,e(1))\) is a pair of
Lemma~\ref{lem:a4-high} and \(s_-\) is a rate of Lemma~\ref{lem:a4-low} under those rules. Steps 2 to 5 establish these
two conditions.

\emph{Step 2: the quiet market.} Let \(\rho\DT\le2k\), and take \(s_+=s_-=0\) and entry \(\rho\) after every order.
Every posterior equals \(1/2<\tau\). The pair \((0,\rho)\) is the pair (N) of Lemma~\ref{lem:a4-high}, because
\(\rho\DT(1-\pi)=\rho\DT/2\le k\), and the rate \(0\) is the rate (N) of Lemma~\ref{lem:a4-low}, because
\(\rho\DT\pi=\rho\DT/2\le k\). Neither condition involves a tie at the threshold, so the quiet market is an
equilibrium under either challenger rule.

\emph{Step 3: the low-class rate.} Let \(\rho\DT>2k\), so that \(t>0\). By Step 1, \(\rho\DT m\ge k\) holds exactly
when \(b\,(\rho\DT/k-2)\ge\lambda\). At the prior \(1/2\) the rate \(\hat s_-\) of Lemma~\ref{lem:a4-low} is
\(b\,(\rho\DT/k-2)\) when \(\rho\DT m<k<\rho\DT/2\), and \(\lambda\) when \(\rho\DT m\ge k\). Hence \(\hat s_-=t\), the one
rate of that lemma, and every tuple in Steps 4 and 5 has \(s_-=t\). Its gross return at entry \(\rho\) satisfies
\begin{equation}\label{eq:o5-sale}
\rho\DT\,\mu_-(t)\ge k ,
\end{equation}
with equality when \(t<\lambda\). Indeed, if \(t=b\,(\rho\DT/k-2)\), then \(2b+t=b\rho\DT/k\) and
\(\mu_-(t)=k/(\rho\DT)\); if \(t=\lambda\), then \(\rho\DT\mu_-(\lambda)=\rho\DT m\ge k\).

\emph{Step 4: a purchase at the rate \(t\) does not recruit.} Let \(\rho\DT>2k\) and \(\mu_+(t)<\tau\). This is the
case whenever \(\tau\ge1\), because \(\mu_+(t)\le M<1\). Take \(s_+=t\) and \(e(1)=\rho\). If \(t<\lambda\), then
\((t,\rho)\) is the pair (Q) of Lemma~\ref{lem:a4-high}: \(0<t<\lambda\), \(\mu_+(t)<\tau\), and
\(\rho\DT[1-\mu_+(t)]=\rho\DT\mu_-(t)=k\) by Step 1 and \eqref{eq:o5-sale}. If \(t=\lambda\), then
\(M=\mu_+(\lambda)<\tau\), so \((\lambda,\rho)\) is the pair (F), and its return condition
\(\rho\DT(1-M)=\rho\DT m\ge k\) holds. The posteriors \(\mu_-(t)<\pi\), \(\mu_0=\pi\) and \(\mu_+(t)\) all lie below
\(\tau\), so the tuple is an equilibrium under either challenger rule, with entry \(\rho\) after every order.

\emph{Step 5: a purchase at the rate \(t\) recruits.} Let \(\rho\DT>2k\) and \(\mu_+(t)\ge\tau\), so that
\(\tau\le\mu_+(t)\le M<1\). Take \(e(1)=1\) and \(s_+=\min\{\lambda,\,b\,(\DT/k-2)\}\). By Step 1 and
\eqref{eq:o5-sale},
\begin{equation}\label{eq:o5-buy}
\DT[1-\mu_+(t)]=\DT\mu_-(t)\ge k/\rho>k .
\end{equation}
Since \(\DT(1-M)=\DT b/(2b+\lambda)\), the inequality \(\DT(1-M)\ge k\) holds exactly when \(b\,(\DT/k-2)\ge\lambda\).

If \(\DT(1-M)\ge k\), then \(s_+=\lambda\ge t\), and \(M=\mu_+(\lambda)\ge\mu_+(t)\ge\tau\). When \(M>\tau\), the pair
(F) of Lemma~\ref{lem:a4-high} has \(e(1)=1\); when \(M=\tau\), either challenger rule admits \(e(1)=1\). The return
condition \(e(1)\DT(1-M)=\DT(1-M)\ge k\) holds, so \((\lambda,1)\) is the pair (F).

If \(\DT(1-M)<k\), then \(s_+=b\,(\DT/k-2)<\lambda\), and \(t<\lambda\), because \(t=\lambda\) would give
\(\DT(1-M)=\DT[1-\mu_+(t)]>k\) by \eqref{eq:o5-buy}. The return \(\DT[1-\mu_+(s)]=\DT b/(2b+s)\) is continuous and
strictly decreasing in \(s\). It exceeds \(k\) at \(s=t\) by \eqref{eq:o5-buy}, and it equals \(k\) exactly at
\(s=s_+\). Hence \(t<s_+<\lambda\). Since \(\mu_+\) rises strictly, \(\mu_+(s_+)>\mu_+(t)\ge\tau\), so
Lemma~\ref{lem:prices}(d) gives entry one after a purchase under either challenger rule, and \((s_+,1)\) is the pair
(R) of Lemma~\ref{lem:a4-high}.

In both cases \(s_+\ge t=s_-\). The posteriors after a sale and after the zero order are at most \(1/2<\tau\), so
their entry is \(\rho\) under either challenger rule, and Step 1 shows that the tuple is an equilibrium.
\end{proof}
